\chapter{1909:Wilhelm Ostwald}
\label{chap:chemistry1909}

The Nobel Prize in Chemistry for the year 1909 was awarded to Wilhelm Ostwald.

\textbf{Motivation:}
in recognition of his work on catalysis and for his investigations into the fundamental principles governing chemical equilibria and rates of reaction

\section{Wilhelm Ostwald}

\subsection*{Early Life and Education}

Wilhelm Ostwald was born on 1853-09-02 in Riga, Russian Empire (now Latvia).
Wilhelm Friedrich Ostwald was a Baltic German chemist and philosopher. Ostwald is credited with being one of the founders of the field of physical chemistry, with Jacobus Henricus van 't Hoff, Walther Nernst and Svante Arrhenius.
He received the Nobel Prize in Chemistry in 1909 for his scientific contributions to the fields of catalysis, chemical equilibria and reaction velocities.

\subsection*{Career and Major Research}

Ostwald was professor of physical chemistry at Leipzig, where he founded the Zeitschrift für physikalische Chemie with van 't Hoff and Arrhenius. He studied catalysis, chemical equilibria, and reaction velocities, and he developed the Ostwald process for the oxidation of ammonia to nitric acid. His textbooks established physical chemistry as a distinct discipline.

\subsection*{Nobel Prize-winning Work}

The work for which Wilhelm Ostwald was awarded the Nobel Prize in Chemistry in 1909 was described by the Nobel Committee as: "in recognition of his work on catalysis and for his investigations into the fundamental principles governing chemical equilibria and rates of reaction".

This research fundamentally advanced the field by making groundbreaking contributions that transformed chemical science.

\subsection*{Significance and Impact}

The impact of Wilhelm Ostwald's work extends far beyond the specific achievement recognized by the Nobel Prize.
These contributions continue to influence chemical research and education today.

\subsection*{Later Career and Honors}

Wilhelm Ostwald continued his scientific work until 1932-04-04, passing away in Leipzig, Germany.
He received numerous honors including membership in major scientific academies and societies.

\subsection*{Legacy}

The legacy of Wilhelm Ostwald endures through the lasting impact of his scientific contributions.
Wilhelm Ostwald's work continues to be cited and built upon by researchers across the chemical sciences.

\subsection*{Selected Publications}

Wilhelm Ostwald's major publications include numerous papers in leading journals of the era, detailing the experimental and theoretical work summarized above.

\clearpage

\section*{Chapter Summary}

Wilhelm Ostwald was awarded the prize for his work on catalysis and on the fundamental principles governing chemical equilibria and rates of reaction.
